% Created 2026-03-17 ter 15:16
% Intended LaTeX compiler: lualatex
\documentclass[12pt]{article}
\usepackage[a4paper,left=3cm,top=3cm,right=2cm,bottom=2cm]{geometry}
\setcounter{secnumdepth}{4}
\setcounter{tocdepth}{2}
\usepackage{graphicx}
\usepackage{longtable}
\usepackage{wrapfig}
\usepackage{rotating}
\usepackage[normalem]{ulem}
\usepackage{amsmath}
\usepackage{amssymb}
\usepackage{capt-of}
\usepackage{hyperref}
\usepackage{fgv-header}
\author{Gustavo M. Mendes de Tarso}
\date{\today}
\title{}
\hypersetup{
 pdfauthor={Gustavo M. Mendes de Tarso},
 pdftitle={},
 pdfkeywords={},
 pdfsubject={},
 pdfcreator={Emacs 28.2 (Org mode 9.5.5)}, 
 pdflang={English}}
\begin{document}

\begingroup
\linespread{1}\selectfont
\begin{tcolorbox}[title=Ficha Técnica,
  colback=gray!5,colframe=gray!40,boxrule=0.4pt,sharp corners]
\begin{tblr}{rowsep=1pt,stretch=1, rows={t}, colspec={Q[l,2.8cm] X[l]}}
\textbf{Curso}                   & Mestrado de Políticas Públicas e Governo \\
\textbf{Turma}                   & T-01 \\
\textbf{Pólo}                    & Brasília \\
\textbf{Disciplina}              & Teorias de Administração Pública \\
\textbf{Professor}               & Bernardo Buta \\
\textbf{Trabalho}                & Aplicação do protocolo PRISMA 2020 a artigo empírico \\
\textbf{Aluno(s)}                & Gustavo M. Mendes de Tarso \\
\textbf{Título do trabalho}      & Busca, seleção e síntese de artigo empírico sobre burocracia de nível de rua, tecnologia e inteligência artificial \\
\end{tblr}
\end{tcolorbox}
\endgroup

\FGVNewPage

\section{Introdução}
\label{sec:orge7ff8bf}

Este trabalho apresenta, em formato alinhado ao \textbf{PRISMA 2020}, o fluxo de identificação, triagem, elegibilidade e inclusão de um \textbf{artigo acadêmico empírico} relacionado ao tema da disciplina \textbf{Teorias de Administração Pública}, com ênfase em sua articulação com \textbf{tecnologia} e \textbf{inteligência artificial}. O recorte adotado procurou preservar o mesmo ponto de partida metodológico exposto em aula: busca no campo da \textbf{Administração Pública} a partir de termos localizados no \textbf{título, resumo ou palavras-chave}, com preferência por \textbf{artigos revisados por pares} e publicados em \textbf{inglês}. Sobre esse universo inicial, foram acrescentados descritores relacionados a \textbf{IA}, \textbf{tomada de decisão}, \textbf{burocracia} e \textbf{implementação}.

A escolha desse recorte é coerente com os conteúdos da disciplina por três razões principais. Em primeiro lugar, porque os debates contemporâneos em Administração Pública já incorporam tópicos como \textbf{transformação digital}, \textbf{algoritmos} e \textbf{inteligência artificial} como problemas empíricos relevantes. Em segundo lugar, porque a disciplina enfatiza que a pesquisa em Administração Pública é tipicamente voltada a problemas concretos de governo, gestão e implementação. Em terceiro lugar, porque a própria estrutura analítica mobilizada em aula valoriza a apresentação clara de \textbf{problema de pesquisa}, \textbf{objetivo}, \textbf{método} e \textbf{achados}, exatamente o formato adotado neste trabalho.

Como esta atividade foi desenvolvida fora de uma interface autenticada de base indexadora, a busca foi \textbf{operacionalizada com a mesma lógica apresentada na aula} e a seleção final foi verificada a partir de \textbf{metadados públicos, páginas de periódicos e versões abertas dos artigos}. O objetivo não foi produzir uma revisão exaustiva, mas demonstrar, de maneira transparente, um processo de busca e seleção compatível com a declaração PRISMA 2020.

Ao final do processo, foi selecionado o artigo \textbf{“Just like I thought”: Street-level bureaucrats trust AI recommendations if they confirm their professional judgment} (Selten, Robeer e Grimmelikhuijsen, 2023), por ser o estudo que melhor conecta um dos eixos centrais da disciplina — a \textbf{burocracia de nível de rua} — com o uso de \textbf{inteligência artificial} na decisão pública.

\section{Estratégia de busca}
\label{sec:org3737226}

\subsection{Pergunta orientadora}
\label{sec:org8ad7521}

\begin{quote}
Qual artigo empírico, aderente ao campo da administração pública, melhor articula os temas da disciplina com tecnologia e inteligência artificial, especialmente nos eixos burocracia, discricionariedade, implementação e tomada de decisão?
\end{quote}

\subsection{Base e lógica de busca}
\label{sec:org807b51d}

A busca partiu do mesmo padrão metodológico exibido no material da aula:
\begin{itemize}
\item \textbf{Scopus} como referência do corpus de Administração Pública;
\item ocorrência do termo de interesse em \textbf{título, resumo ou palavras-chave};
\item preferência por \textbf{artigos revisados por pares} e \textbf{publicados em inglês}.
\end{itemize}

A partir desse núcleo, foram acrescidos descritores relacionados a inteligência artificial e tecnologia, seguidos de triagem manual para retenção apenas de \textbf{artigos empíricos} com aderência efetiva ao debate da disciplina.

\subsection{Query de trabalho (formato compatível com a lógica da Scopus)}
\label{sec:org7cf57ad}

\begin{quote}
TITLE-ABS-KEY("public administration")
AND TITLE-ABS-KEY("artificial intelligence" OR algorithm* OR "digital transformation")
AND TITLE-ABS-KEY("street-level bureaucr*" OR bureaucr* OR "decision making" OR governance)
AND NOT TITLE-ABS-KEY(review OR "systematic review" OR "literature review")
\end{quote}

\subsection{Fontes operacionais consultadas nesta sessão}
\label{sec:org3479cc7}
\begin{itemize}
\item Material da disciplina (\textbf{slides} e textos-base)
\item Páginas abertas de periódicos científicos
\item Repositórios institucionais com metadados e/ou PDF
\item Registros DOI e páginas de metadados acadêmicos
\end{itemize}

\subsection{Critérios de inclusão}
\label{sec:orgf0baf58}
\begin{itemize}
\item artigo de periódico;
\item estudo empírico;
\item aderência ao campo da administração pública;
\item presença explícita de tecnologia e/ou inteligência artificial no problema de pesquisa;
\item conexão substantiva com ao menos um eixo da disciplina, como burocracia, discricionariedade, implementação, governança ou tomada de decisão;
\item disponibilidade de resumo e de informações metodológicas suficientes para avaliação, preferencialmente com acesso ao texto completo.
\end{itemize}

\subsection{Critérios de exclusão}
\label{sec:org072cb97}
\begin{itemize}
\item revisões de literatura, artigos ensaísticos ou agendas de pesquisa;
\item preprints não revisados por pares;
\item estudos com foco predominantemente técnico-computacional, sem centralidade para o debate em administração pública;
\item estudos empíricos sobre adoção tecnológica cujo vínculo com os eixos teóricos centrais da disciplina fosse apenas periférico.
\end{itemize}

\section{PRISMA 2020}
\label{sec:org583665e}

\subsection{Fluxo textual dos estudos}
\label{sec:org64238bb}

\subsubsection{Identificação}
\label{sec:org034e22f}
\begin{itemize}
\item Registros identificados por busca focalizada: \textbf{n = 8}
\begin{itemize}
\item Registros sobre IA, decisão pública e burocracia/organizações públicas: \textbf{n = 8}
\end{itemize}

\item Registros removidos antes da triagem: \textbf{n = 1}
\begin{itemize}
\item Duplicata do artigo selecionado em página de periódico e repositório institucional: \textbf{n = 1}
\item Remoção por outros motivos: \textbf{n = 0}
\end{itemize}
\end{itemize}

\subsubsection{Triagem}
\label{sec:org03fb0ca}
\begin{itemize}
\item Registros triados por título, resumo e metadados: \textbf{n = 7}
\item Registros excluídos na triagem: \textbf{n = 3}
\begin{itemize}
\item Revisão de literatura, scoping review ou agenda de pesquisa: \textbf{n = 2}
\item Preprint não revisado por pares: \textbf{n = 1}
\end{itemize}
\end{itemize}

\subsubsection{Elegibilidade}
\label{sec:orgdaef0e9}
\begin{itemize}
\item Relatórios buscados para recuperação em texto completo: \textbf{n = 4}
\item Relatórios não recuperados: \textbf{n = 0}

\item Relatórios avaliados para elegibilidade: \textbf{n = 4}
\item Relatórios excluídos após leitura do texto completo: \textbf{n = 3}
\begin{itemize}
\item Foco principal em adoção organizacional de IA, com ligação apenas indireta à burocracia de nível de rua e à discricionariedade: \textbf{n = 2}
\item Ênfase em accountability subjetiva mediada por IA, menos aderente ao recorte priorizado de decisão e implementação: \textbf{n = 1}
\end{itemize}
\end{itemize}

\subsubsection{Inclusão}
\label{sec:orgc9b47a2}
\begin{itemize}
\item Estudos incluídos na síntese qualitativa: \textbf{n = 1}
\end{itemize}

\subsection{Tabela-resumo do fluxo PRISMA 2020}
\label{sec:orgd639311}

\begin{table}[htbp]
\centering
\scriptsize
\setlength{\tabcolsep}{4pt}
\renewcommand{\arraystretch}{1.15}
\begin{tabularx}{\textwidth}{>{\raggedright\arraybackslash}X >{\raggedleft\arraybackslash}m{0.18\textwidth}}
\toprule
Etapa PRISMA 2020 & n \\
\midrule
Registros identificados & 8 \\
Registros removidos antes da triagem & 1 \\
Registros triados & 7 \\
Registros excluídos na triagem & 3 \\
Relatórios buscados para recuperação & 4 \\
Relatórios não recuperados & 0 \\
Relatórios avaliados para elegibilidade & 4 \\
Relatórios excluídos após leitura completa & 3 \\
Estudos incluídos na síntese qualitativa & 1 \\
\bottomrule
\end{tabularx}
\normalsize
\end{table}

\subsection{Estudos classificados no fluxo}
\label{sec:orgb72bd66}

\begin{table}[htbp]
\centering
\scriptsize
\setlength{\tabcolsep}{3pt}
\renewcommand{\arraystretch}{1.15}
\begin{tabularx}{\textwidth}{>{\raggedright\arraybackslash}X >{\raggedright\arraybackslash}m{0.18\textwidth} >{\raggedright\arraybackslash}m{0.14\textwidth} >{\raggedright\arraybackslash}m{0.22\textwidth}}
\toprule
Estudo & Situação no fluxo & Tipo & Motivo \\
\midrule
Selten, Robeer e Grimmelikhuijsen (2023), \textit{“Just like I thought”: Street-level bureaucrats trust AI recommendations if they confirm their professional judgment} & Incluído & Artigo empírico & Melhor articulação entre burocracia de nível de rua, discricionariedade e IA \\
Deng e Sun (2025), \textit{How does the usage of artificial intelligence affect felt administrative accountability of street-level bureaucrats?} & Excluído na elegibilidade & Artigo empírico & Foco principal em accountability percebida, menos aderente ao recorte de decisão e implementação \\
Ribeiro e Segatto (2025), \textit{Inteligência artificial nas organizações públicas brasileiras: heterogeneidades e capacidades em tecnologia da informação} & Excluído na elegibilidade & Artigo empírico & Estudo organizacional relevante, mas com vínculo indireto ao eixo de burocracia de nível de rua \\
Omonov e Ahn (2025), \textit{Towards Smart Public Administration: A TOE-Based Empirical Study of AI Chatbot Adoption in a Transitioning Government Context} & Excluído na elegibilidade & Artigo empírico & Ênfase em adoção de chatbot e framework TOE, com menor diálogo com os eixos teóricos centrais da disciplina \\
Kaun, Larsson e colaboradores (2024), \textit{Automating public administration: citizens' attitudes towards automated decision-making across Estonia, Sweden, and Germany} & Excluído na triagem & Artigo empírico & Foco em atitudes dos cidadãos, e não na burocracia, organização ou implementação \\
Gillingham, Morley e Floridi (2024), \textit{The Effects of AI on Street-Level Bureaucracy: A Scoping Review} & Excluído na triagem & Revisão de literatura & Não é artigo empírico \\
Johannesson e Arnesen (2025), \textit{How Using Artificial Intelligence in Public Administration Shifts Citizens' Expectations of Street-Level Bureaucracy} & Excluído na triagem & Preprint & Não revisado por pares \\
\bottomrule
\end{tabularx}
\normalsize
\end{table}

\clearpage
\clearpage
\subsection{Representação esquemática em PlantUML}
\label{sec:org2f52f52}

\FGVAutoSVG
  {prisma_empirico_ia_teorias_adm_publica_fluxo}
  {Fluxograma PRISMA 2020 da seleção do artigo empírico.}
  {fig:prisma2020_fluxo}

\section{Resumo do artigo selecionado}
\label{sec:org74649b8}

\subsection{Referência do estudo}
\label{sec:orgfe084f4}

Selten, Friso; Robeer, Marcel; Grimmelikhuijsen, Stephan. \textbf{“Just like I thought”: Street-level bureaucrats trust AI recommendations if they confirm their professional judgment}. \emph{Public Administration Review}, v. 83, n. 2, p. 263-278, 2023.

\subsection{Problema e objetivo de pesquisa}
\label{sec:orgdb1b6df}

O problema central do artigo é diretamente compatível com um dos núcleos da disciplina: \textbf{como burocratas de nível de rua interagem com tecnologias de inteligência artificial quando essas tecnologias passam a participar da tomada de decisão pública}. Mais especificamente, o estudo procura responder o que acontece quando uma recomendação produzida por IA \textbf{confirma} ou \textbf{contradiz} o julgamento profissional intuitivo de um burocrata de nível de rua.

O objetivo do artigo é duplo. O primeiro é investigar de que forma recomendações de IA afetam a \textbf{decisão} de burocratas de nível de rua. O segundo é examinar se a chamada \textbf{explainable AI} aumenta a confiança dos agentes públicos nessas recomendações. Em termos teóricos, o estudo confronta duas possibilidades: de um lado, a \textbf{automation bias}, isto é, a tendência de seguir a máquina; de outro, a \textbf{confirmation bias}, isto é, a tendência de confiar mais na IA quando ela apenas reforça o que o agente já acreditava.

Essa formulação faz uma ponte clara entre os temas da aula, porque recoloca em chave contemporânea debates clássicos sobre \textbf{discricionariedade}, \textbf{implementação}, \textbf{autonomia decisória} e \textbf{controle} no nível da burocracia de rua. Em vez de perguntar apenas como o burocrata interpreta normas, o artigo pergunta como ele reage quando uma tecnologia algorítmica entra no processo decisório.

\subsection{Método utilizado}
\label{sec:orgcc0d1e6}

O artigo adota um \textbf{desenho experimental} com forte aderência empírica ao contexto profissional estudado. Os autores aplicaram um \textbf{experimento fatorial 2 × 2, com medidas repetidas}, utilizando um sistema simulado de \textbf{predictive policing} concebido para ser realista no contexto da polícia holandesa.

A amostra foi composta por \textbf{124 policiais de nível de rua} nos Países Baixos. Cada participante respondeu a \textbf{três vinhetas} semelhantes, resultando em \textbf{294 observações} no total. Os tratamentos combinaram duas dimensões:
\begin{itemize}
\item recomendação de IA \textbf{congruente} ou \textbf{incongruente} com o julgamento profissional intuitivo do policial;
\item recomendação \textbf{explicada} ou \textbf{não explicada} pela IA.
\end{itemize}

Com isso, o estudo pôde testar, de maneira controlada, se os policiais \textbf{confiam}, \textbf{aceitam} e \textbf{seguem} mais uma recomendação algorítmica quando ela está alinhada ao seu conhecimento profissional prévio e se a explicação do sistema altera esse comportamento.

Do ponto de vista das diretrizes de pesquisa discutidas em aula, trata-se sobretudo de uma \textbf{pesquisa explicativa}, pois busca identificar o efeito de uma condição causal — o tipo de recomendação da IA — sobre crenças e comportamentos de agentes públicos.

\subsection{Principais achados}
\label{sec:org1a2f4c2}

Os resultados do estudo mostram que os policiais \textbf{não seguem cegamente} as recomendações da IA. Ao contrário, eles tendem a \textbf{confiar} e \textbf{seguir} essas recomendações quando elas \textbf{confirmam} seu julgamento profissional anterior. Em outras palavras, o padrão dominante observado é mais compatível com \textbf{viés de confirmação} do que com \textbf{viés de automação}.

O segundo achado relevante é que a \textbf{IA explicável} não produziu aumento estatisticamente significativo de confiança nas recomendações algorítmicas. Isso significa que, ao menos neste contexto, oferecer explicações adicionais sobre a lógica da recomendação não foi suficiente para superar o peso do julgamento profissional já formado pelo agente.

O terceiro achado é normativamente ambivalente. Por um lado, a persistência da discricionariedade humana pode funcionar como mecanismo de correção de recomendações algorítmicas enviesadas, injustas ou falhas. Por outro lado, o mesmo padrão limita a expectativa de que sistemas de IA, mesmo quando concebidos para promover maior justiça ou neutralidade, consigam corrigir automaticamente os vieses humanos. O estudo mostra, portanto, que a interação entre burocrata e algoritmo é mais complexa do que uma simples substituição da decisão humana pela decisão automatizada.

\subsection{Por que este artigo foi o selecionado}
\label{sec:orgc805ddd}

Entre os estudos empíricos encontrados, este foi o mais aderente ao conteúdo da disciplina por quatro razões principais. Primeiro, porque se ancora diretamente na literatura de \textbf{burocracia de nível de rua}. Segundo, porque trata de \textbf{discricionariedade} e \textbf{tomada de decisão}, dois temas centrais do campo. Terceiro, porque examina um problema contemporâneo da administração pública — a incorporação de IA em decisões operacionais — sem abandonar a linguagem teórica do curso. Quarto, porque oferece um método empírico claro, replicável e facilmente apresentável no formato de \textbf{problema, objetivo, método e achados}.

\section{Considerações finais}
\label{sec:orgd6132f3}

A aplicação do PRISMA 2020 ao recorte aqui adotado mostra que o encontro entre \textbf{teorias da administração pública} e \textbf{inteligência artificial} já produz um conjunto relevante de estudos, mas nem todos dialogam igualmente com os eixos centrais da disciplina. Parte dessa literatura é formada por revisões, agendas de pesquisa ou estudos mais próximos da adoção tecnológica organizacional. Quando o objetivo é encontrar \textbf{um artigo empírico} que articule, de forma mais densa, os temas da aula com IA, o eixo da \textbf{burocracia de nível de rua} se revelou particularmente promissor.

O artigo selecionado evidencia que a introdução de IA na administração pública não elimina o papel do burocrata. Ao contrário, a tecnologia passa a compor um novo ambiente decisório, no qual persistem discricionariedade, julgamento profissional, confiança seletiva e possibilidades de viés. Nesse sentido, a IA não substitui os dilemas clássicos da administração pública; ela os reconfigura.

Por isso, o estudo de Selten, Robeer e Grimmelikhuijsen é particularmente útil para a disciplina. Ele permite atualizar a discussão sobre implementação e discricionariedade com base em um problema contemporâneo, mostrando que a análise de tecnologia pública só ganha densidade explicativa quando é articulada aos conceitos clássicos do campo.

\section{Referências}
\label{sec:org50faae0}

ALON-BARKAT, Sharon; BUSUIOC, Madalina. \textbf{Public Administration Meets Artificial Intelligence: Towards a Meaningful Behavioral Research Agenda on Algorithmic Decision-Making in Government}. Journal of Behavioral Public Administration, 2024.

DENG, Yi; SUN, Yu. \textbf{How does the usage of artificial intelligence affect felt administrative accountability of street-level bureaucrats? The mediating effect of perceived discretion}. Public Management Review, 2025.

GILLINGHAM, Cara; MORLEY, Jessica; FLORIDI, Luciano. \textbf{The Effects of AI on Street-Level Bureaucracy: A Scoping Review}. 2024.

JOHANNESSON, Mikael Poul; ARNESEN, Sveinung. \textbf{How Using Artificial Intelligence in Public Administration Shifts Citizens’ Expectations of Street-Level Bureaucracy}. SocArXiv, 2025.

KAUN, Anne; LARSSON, Anders Olof; et al. \textbf{Automating public administration: citizens' attitudes towards automated decision-making across Estonia, Sweden, and Germany}. Information, Communication \& Society, 2024.

OMONOV, Mansur Samadovich; AHN, Yonghan. \textbf{Towards Smart Public Administration: A TOE-Based Empirical Study of AI Chatbot Adoption in a Transitioning Government Context}. Administrative Sciences, 2025.

RIBEIRO, Manuella Maia; SEGATTO, Catarina Ianni. \textbf{Inteligência artificial nas organizações públicas brasileiras: heterogeneidades e capacidades em tecnologia da informação}. Revista de Administração Pública, v. 59, n. 1, e2024-0066, 2025.

SELTEN, Friso; ROBEER, Marcel; GRIMMELIKHUIJSEN, Stephan. \textbf{“Just like I thought”: Street-level bureaucrats trust AI recommendations if they confirm their professional judgment}. Public Administration Review, v. 83, n. 2, p. 263-278, 2023.
\end{document}